\documentclass[11pt]{article}

\usepackage[margin=1in]{geometry}
\usepackage{amsmath,amssymb,amsthm,mathtools}
\usepackage{enumitem}
\usepackage{hyperref}
\usepackage[nameinlink,capitalize]{cleveref}

\hypersetup{colorlinks=true,linkcolor=blue,citecolor=blue,urlcolor=blue}

\newtheorem{theorem}{Theorem}
\newtheorem{proposition}{Proposition}
\newtheorem{corollary}{Corollary}
\newtheorem{definition}{Definition}
\newtheorem{remark}{Remark}
\newtheorem{warning}{Warning}

\newcommand{\F}{\mathbb F}

\title{The Diagonal Barrier in Kummer Large-Sieve Transfers}
\author{Jared Wilder}
\date{2026-05-13}

\begin{document}
\maketitle

\begin{abstract}
Paper 4 formulated a rank-two Hecke large-sieve bridge for arbitrary coefficient weights on Kummer vectors.  This note proves that such a bridge is false as stated.  The finite vector space \(\F_\ell^2\) imposes an exact Plancherel identity with an unavoidable diagonal and fiber-sum term.  A single localized coefficient already contributes \(\ell^2-1\), contradicting any bound of size \(\ell^2(\log Q)^{-A}\) for arbitrary unit \(L^2\) weights.  The correct finite object is the vector-fiber variance, not a universal arbitrary-coefficient large sieve.  Projective slope collision control does not imply vector-fiber control.  The Bernoulli-Kummer smoothing theorem supplies a distributional version for the obstruction law, but Kummer-BV/BVK still require a separate sieve-transfer estimate for the actual von Mangoldt or sieve weights.
\end{abstract}

\section{Setup}

Let
\[
G=\F_\ell^2.
\]
For each prime \(q\in\mathcal Q\), let
\[
x_q=(\log_q2,\log_q3)\in G.
\]
For complex weights \(\omega_q\), define the Fourier sums
\[
S_A(\omega)=\sum_{q\in\mathcal Q}\omega_q e_\ell(A\cdot x_q),
\qquad A\in G,
\]
and the vector-fiber sums
\[
\Omega(v)=\sum_{q:x_q=v}\omega_q,
\qquad v\in G.
\]

\section{Exact identity}

\begin{theorem}[Exact Kummer large-sieve identity]\label{T:exact}
For arbitrary complex weights \(\omega_q\),
\[
\sum_{A\ne0}|S_A(\omega)|^2
=
\ell^2\sum_{v\in G}|\Omega(v)|^2
-
\left|\sum_{q\in\mathcal Q}\omega_q\right|^2.
\]
\end{theorem}

\begin{proof}
Since \(S_A(\omega)=\sum_{v\in G}\Omega(v)e_\ell(A\cdot v)\), Plancherel on \(G\) gives
\[
\sum_{A\in G}|S_A(\omega)|^2
=
\ell^2\sum_{v\in G}|\Omega(v)|^2.
\]
The \(A=0\) term is
\[
|S_0(\omega)|^2
=
\left|\sum_q\omega_q\right|^2.
\]
Subtracting it proves the identity.
\end{proof}

\begin{corollary}[Diagonal counterexample]\label{C:single}
The bound
\[
\sum_{A\ne0}|S_A(\omega)|^2
\ll_A
\ell^2(\log Q)^{-A}
\]
cannot hold for all weights satisfying \(\sum_q|\omega_q|^2\le1\).
\end{corollary}

\begin{proof}
Choose one prime \(q_0\) and set \(\omega_{q_0}=1\), \(\omega_q=0\) for \(q\ne q_0\).  Then
\[
|S_A(\omega)|^2=1
\qquad(A\ne0),
\]
so
\[
\sum_{A\ne0}|S_A(\omega)|^2=\ell^2-1.
\]
For fixed \(\ell\), the proposed right-hand side tends to \(0\) relative to \(\ell^2\) as \(Q\to\infty\).  In the high-\(\ell\) range it is still smaller by an arbitrary logarithmic power.  The stated arbitrary-coefficient saving is therefore false.
\end{proof}

\begin{corollary}[Fiber multiplicity barrier]\label{C:fiber}
If \(M\) indices \(q\) share the same vector \(x_q=v_0\), and \(\omega_q=M^{-1/2}\) on those indices and \(0\) elsewhere, then \(\sum_q|\omega_q|^2=1\) but
\[
\sum_{A\ne0}|S_A(\omega)|^2
=
M(\ell^2-1).
\]
Thus the operator norm is controlled by vector-fiber multiplicity, not only by the ambient dimension.
\end{corollary}

\begin{proof}
Here \(\Omega(v_0)=M^{1/2}\), all other fibers vanish, and \(\sum_q\omega_q=M^{1/2}\).  Substitution into \cref{T:exact} gives
\[
\ell^2M-M=M(\ell^2-1).
\]
\end{proof}

\section{Distributional version}

The identity becomes useful when the weights define a probability distribution rather than arbitrary localized coefficients.

\begin{proposition}[Probability-weight Plancherel]\label{P:prob}
Let
\[
\mu=\sum_{q\in\mathcal Q}a_q\delta_{x_q},
\qquad a_q\ge0,\qquad \sum_q a_q=1.
\]
Then
\[
\sum_{A\ne0}|\widehat\mu(A)|^2
=
\ell^2\|\mu-u\|_2^2,
\]
where \(u\) is uniform on \(G\).
\end{proposition}

\begin{proof}
Apply \cref{T:exact} with \(\omega_q=a_q\).  The fiber sums are the probabilities \(\mu(v)\), and the \(A=0\) term is \(1\).  Since
\[
\|\mu-u\|_2^2
=
\sum_v\mu(v)^2-\ell^{-2},
\]
we get
\[
\ell^2\sum_v\mu(v)^2-1
=
\ell^2\|\mu-u\|_2^2.
\]
\end{proof}

\begin{corollary}[Bernoulli-Kummer distributional large sieve]\label{C:bernoulli}
Let \(\nu\) be the Bernoulli-Kummer obstruction law from Paper 12.  If
\[
W_\ell(\mathcal Q)=\min_{A\ne0}\sum_{q:A\cdot x_q\ne0}p_q(1-p_q),
\]
then
\[
\sum_{A\ne0}|\widehat\nu(A)|^2
\le
(\ell^2-1)
\exp\left(-\frac{16}{\ell^2}W_\ell(\mathcal Q)\right).
\]
In particular, if \(16W_\ell(\mathcal Q)/\ell^2\ge A\log\log Q+O(\log\ell)\), then
\[
\sum_{A\ne0}|\widehat\nu(A)|^2
\ll_A(\log Q)^{-A}.
\]
\end{corollary}

\begin{proof}
Paper 12 proves
\[
\|\nu-u\|_2^2
\le
(1-\ell^{-2})
\exp\left(-\frac{16}{\ell^2}W_\ell(\mathcal Q)\right).
\]
Multiply by \(\ell^2\) and use \cref{P:prob}.
\end{proof}

\section{Projective slope control is not vector-fiber control}

The projective Kummer collision graph records the class
\[
\sigma_\ell(q)=[x_q]\in\mathbb P^1(\F_\ell)
\]
when \(x_q\ne0\).  This forgets the radial scalar in \(x_q\).  Additive characters \(e_\ell(A\cdot x_q)\), however, see the full vector.

\begin{proposition}[Projective-to-vector gap]\label{P:projective-gap}
Small projective slope collision does not imply small vector-fiber variance for arbitrary coefficients.
\end{proposition}

\begin{proof}
Choose one nonzero vector on each projective line in \(\F_\ell^2\), and put equal probability mass on these \(\ell+1\) representatives.  The projective slope distribution is perfectly uniform.  But the vector distribution has support size only \(\ell+1\), so
\[
\|\mu-u\|_2^2
=
\frac1{\ell+1}-\frac1{\ell^2},
\]
which is much larger than a logarithmic power saving target in the high-\(\ell\) range.  By \cref{P:prob}, the corresponding nontrivial character energy is
\[
\ell^2\left(\frac1{\ell+1}-\frac1{\ell^2}\right)
\asymp \ell.
\]
Thus slope equidistribution alone does not give vector Fourier saving.  Arbitrary coefficients localized on a single vector fiber give the stronger obstruction in \cref{C:single,C:fiber}.
\end{proof}

\section{Corrected sieve-transfer status}

The finite correction is:

\begin{itemize}[leftmargin=2em]
\item There is no arbitrary-coefficient Kummer large-sieve saving of size \(\ell^2(\log Q)^{-A}\) under only \(\sum|\omega_q|^2\le1\).
\item The exact finite quantity is the vector-fiber energy \(\sum_v|\Omega(v)|^2\).
\item Projective KRC is a slope-statistic; it does not control radial vector fibers.
\item Bernoulli-Kummer smoothing supplies a distributional large-sieve statement for the obstruction probability law.
\item Kummer-BV and BVK still require a separate analytic transfer statement for the actual von Mangoldt or sieve weights.  That transfer must control vector-fiber sums for those weights, or state an equivalent Fourier variance estimate directly.
\end{itemize}

\begin{warning}[What this does not prove]
This note does not prove Kummer-BV, BVK, or Erd\H{o}s--Graham \#203.  It proves that one proposed arbitrary-coefficient bridge is false and replaces it with the exact finite identity that any valid sieve-transfer theorem must satisfy.
\end{warning}

\end{document}
